\section{The Five Primitives}
\label{sec:primitives}

Algorithm~\ref{alg:harness} gives a harness as pseudocode, each line labelled with its primitive. Every harness in \S\ref{sec:deconstruction} is an elaboration of it. For each primitive we give its \emph{mechanism} (what it does on the path of a request), its \emph{negative definition} (what it never does, and which primitive does), and its \emph{design space} (the decisions a harness makes within it). The negative definitions make the five mutually exclusive: each piece of a harness has one home, assigned by its function rather than by its name or its location in the code.

\begin{algorithm}[t]
\caption{A harness. Lines are labelled with the primitive they implement.}
\label{alg:harness}
\small
\begin{algorithmic}[1]
\Procedure{ControlFlow}{\,}
  \State $\textit{conversation} \gets [\,]$
  \Loop \Comment{outer loop: one request at a time}
    \State $\textit{conversation}.\textsc{append}(\textsc{User}(\Call{Input}{\,}))$ \Comment{\prim{Input}: from a person}
    \Loop \Comment{inner loop: one request, carried through}
      \State $\textit{ctx} \gets \Call{Context}{\textit{conversation}}$ \Comment{\prim{Context}}
      \State $\textit{response} \gets \Call{ModelInterface}{\textit{ctx}}$ \Comment{\prim{Model interface}}
      \State $\textit{conversation}.\textsc{append}(\textsc{Assistant}(\textit{response}))$
      \State $\textit{results} \gets \Call{Output}{\textit{response}}$ \Comment{\prim{Output}: show words, run actions}
      \If{$\textit{results} = \emptyset$} \textbf{break} \Comment{nothing ran: hand back}
      \EndIf
      \State $\textit{conversation}.\textsc{append}(\textsc{User}(\textit{results}))$ \Comment{\prim{Input}: from the world}
    \EndLoop
  \EndLoop
\EndProcedure
\end{algorithmic}
\end{algorithm}

\subsection{\prim{Input}: capture what comes in}
\textbf{Mechanism.} Capture what enters the agent and hand it to \prim{Context}. There are two sources: a person (a message, a file, an interrupt, an answer to a question the model asked) and the world (the result of an action, an event, a scheduled trigger).
\textbf{Never.} \prim{Input} never decides what the model is shown; that is \prim{Context}. It never runs anything; that is \prim{Output}. An answer a person gives that the \emph{harness} consumes, such as an approval, is not input but a \prim{Control flow} decision; an answer the \emph{model} consumes, returned as a tool result, is input.
\textbf{Design space.} The source: a terminal, an editor, a chat channel, a command-line argument, an issue tracker, a webhook, a message queue. The kind: text, images, files, audio, the results of actions, interrupts, and mid-turn steering. The second source, results from the world, is what makes acting worthwhile: an action's result is new perception.

\subsection{\prim{Context}: assemble what the model is shown}
\textbf{Mechanism.} Assemble, for one call, everything the model is shown: instructions, the tools' descriptions, the conversation so far, and anything retrieved from memory or the environment; and fit it to the window, by truncation, compaction or selection. \prim{Context} is the agent's working memory, rebuilt on every call.
\textbf{Never.} \prim{Context} never calls the model; that is \prim{Model interface}. It never acts on the world; that is \prim{Output}. Reading a memory store is \prim{Context}; writing one, if the model requests it, is an action and so \prim{Output}. Code that decides what the model sees is \prim{Context} even when it sits inside a tool: a tool that attaches project instructions to a file it reads, or condenses a web page before returning it, is doing \prim{Context}'s work.
\textbf{Design space.} What goes in the instructions (identity, rules, the agent's own capabilities, project files such as \texttt{AGENTS.md}); how tools are described; how long-term memory comes in (preloaded, retrieved on demand, or searched by the agent through its own tools); the roles each piece is placed in, which carry authority; and what happens when the window fills: drop old turns, summarize them, or include only what the moment needs.

\subsection{\prim{Model interface}: reach the model}
\textbf{Mechanism.} Send the assembled context to the model and bring back its response in a form the rest of the harness can read: words, and structured requests to act. This includes the client, the wire format, the translation between the harness's message format and the model's chat template, streaming, retrying a failed request, and switching to a backup model when the request could not be served.
\textbf{Never.} \prim{Model interface} never decides what the request contains; that is \prim{Context}. It never acts on the response; that is \prim{Output}. It receives a stream; rendering each fragment as it arrives is \prim{Output}. A change of model because the request could not be served (unavailable, overloaded, refused) is \prim{Model interface}; a change of model keyed to the work (planning, a cheap model for a side task) is \prim{Control flow}.
\textbf{Design space.} A hosted API or local inference; one provider or many, directly or through a translation layer; native tool-call formats or tool calls parsed from text; streaming or whole responses; retry policy and backup models. \prim{Model interface} is the only primitive that knows which model is inside, which is why changing the model means changing it alone.

\subsection{\prim{Output}: carry the response out}
\textbf{Mechanism.} Handle what the model returned: show its words to a person, and carry out each requested action, collecting its result. The tool set, the executor that runs an action, and the environment it runs in are \prim{Output}.
\textbf{Never.} \prim{Output} never decides whether to go again; that is \prim{Control flow}. It never assembles the next request; that is \prim{Context}. The model only ever \emph{requests} an action; \prim{Output} decides whether and how it happens, which is where sandboxes and never-lists live. Several tool calls from one response run concurrently: \prim{Output}. Several \emph{model} calls run concurrently: \prim{Control flow}. A program the model writes at run time that calls tools without calling the model is \prim{Output}; one that sequences model calls is \prim{Control flow}.
\textbf{Design space.} The tool set, from one general tool (a shell) to dozens of specific ones (read, write, edit, search, browse, computer use); the action format, from structured tool calls to code the model writes and the harness executes \citep{wang2024codeact} to edit formats such as search-and-replace blocks; the execution environment, from the host machine to a container or a remote sandbox; and the presentation of results, including truncation and error messages, which \citet{yang2024sweagent} show affect performance.

\subsection{\prim{Control flow}: decide what happens next}
\textbf{Mechanism.} Define the shape of the loop: send results back and go again, hand back to the person, or stop. Step and token budgets, approval gates the harness consumes, delegation to sub-agents and the sequencing of several agents, and routing to a model keyed to the work are all \prim{Control flow}.
\textbf{Never.} \prim{Control flow} never captures input, assembles context, calls the model or runs a tool. It sequences the other four. A workflow in which code fixes the order of model calls \citep{anthropic2024effective} and an agent in which the model directs its own process differ here and only here.
\textbf{Design space.} A single open loop that runs until the model stops requesting actions; a fixed pipeline; an explicit graph; sub-agents with contexts of their own; planning or approval modes that change what is allowed; and stopping rules: step limits, budgets, interrupts, loop detection.

\subsection{Necessity}
Each primitive answers a question no other answers: what comes in, what the model is shown, how the model is reached, what is done with what comes out, and what happens next. Appendix~\ref{app:minimal} gives the five minimal harnesses from which the definitions were drawn, built one primitive at a time around a fixed model and task. Each version fails in a specific way without the next primitive: with only \prim{Model interface}, the model answers but cannot act and nobody can ask it anything else; with \prim{Input}, anyone can ask and it still cannot act; with \prim{Output}, it acts once and the result goes nowhere; with \prim{Control flow}, it finishes the task but does not know what it is or where it is, and spends actions finding out; with \prim{Context}, it knows. This shows that each primitive removes a failure the others do not. It does not show that no sixth primitive could ever be needed; \S\ref{sec:deconstruction} tests that against existing harnesses, and \S\ref{sec:limitations} says what would refute it.
